\documentclass[11pt,a4paper]{amsart}
\usepackage[T1]{fontenc}
\usepackage{amsmath,amssymb}
\usepackage{microtype}
\usepackage[margin=1in]{geometry}
\usepackage[unicode,hidelinks]{hyperref}
\emergencystretch=1em

\newtheorem{theorem}{Theorem}[section]
\newtheorem{lemma}[theorem]{Lemma}
\newtheorem{proposition}[theorem]{Proposition}
\newtheorem{corollary}[theorem]{Corollary}
\theoremstyle{definition}
\newtheorem{definition}[theorem]{Definition}
\theoremstyle{remark}
\newtheorem{remark}[theorem]{Remark}
\newcommand{\Kminus}{K_7^{-}}

\title[Graphs with no \(K_7^{-}\) minor]{Every graph with no \(K_7^{-}\) minor is six-colourable}
\author{Cavit Erginsoy}
\date{Draft, 24 September 2026}
\hypersetup{pdftitle={Every graph with no K7-minus minor is six-colourable},pdfauthor={Cavit Erginsoy}}

\begin{document}
\begin{abstract}
Let $K_7^{-}$ be the graph obtained from $K_7$ by deleting one edge.
We prove that every graph with no $K_7^{-}$ minor is six-colourable.
The main step is an extremal theorem: every $4$-bilight graph on $n\geq3$
vertices with at least $4n-2$ edges contains a $K_7^{-}$ minor.
In particular, this holds for all five-connected graphs.
The key step strengthens a rooted density theorem of
Dvo\v r\'ak, Norin and Rahman, reducing the number of missing
adjacencies in their construction from two to one.
\end{abstract}
\maketitle

\section{Introduction}
All graphs in this paper are finite and simple. Write $K_7^{-}$ for
$K_7$ with one edge deleted. We prove the following statement, posed as
Conjecture~21 by Norin and Totschnig~\cite{NT2025}.

\begin{theorem}\label{thm:colour}
Every graph with no $K_7^{-}$ minor is six-colourable.
\end{theorem}

Norin and Totschnig~\cite{NT2025} proved the corresponding result when
the excluded minor is $K_7$ with two incident edges deleted.
Dvo\v r\'ak, Norin and Rahman~\cite{DNR2026} proved it for $K_7$ with
two independent edges deleted. Theorem~\ref{thm:colour} strengthens
both conclusions. It does not settle Hadwiger's conjecture for $K_7$.

Our proof uses the rooted density and separation results
of Dvo\v r\'ak, Norin and Rahman. Their critical-graph density bound,
also obtained independently in our earlier work~\cite[Corollary~3]{Erginsoy2026critical},
reduces Theorem~\ref{thm:colour} to an extremal problem. We prove the
stronger density theorem stated as Theorem~\ref{thm:bilight} below,
for the class of $4$-bilight graphs. It gives the following result,
conjectured in~\cite[Conjecture~1.5]{DNR2026}.

\begin{corollary}\label{thm:five-connected}
Every five-connected graph on $n\geq6$ vertices with at least $4n-2$
edges contains a $K_7^{-}$ minor.
\end{corollary}

The rooted construction in~\cite[Theorem~2.9]{DNR2026} gives five
branch sets containing prescribed roots and two additional branch sets.
Of the eleven possible adjacencies involving the latter two sets, it
allows two to be missing. Even when the five rooted branch sets are
pairwise adjacent, this need not give $K_7^{-}$. Our rooted theorem
allows only one missing adjacency when the rooted four-density is at
least two. The difficulty is that an edge contraction may destroy the
separation conditions needed for induction. We show that each such
obstruction has a five-vertex boundary. A rooted-star construction
handles the small configurations. For the rest, we add an auxiliary
clique and apply Mader's atom theorem; its bound on the size of an atom
contradicts a degree count. The auxiliary clique serves only this
argument and is not used to construct a minor.

To prove the unrooted theorem, we must also deal with sides of five-cuts
whose density is at most one. In a minimum counterexample, every five-cut
has exactly one such side. Local contractions and rooted constructions
exclude vertices of degree five; an inclusion-minimal fragment then
excludes all five-cuts. Once six-connectivity is established, edge
contractions force every neighbourhood to have minimum degree at least
four. Minimality gives $e(G)=4v(G)-2$, so some vertex has degree six or seven.
We finish with an adaptation of Norin--Totschnig's degree-six argument
and a short complement classification for degree seven.

Section~\ref{sec:preliminaries} gives the definitions and external
results. The rooted theorem and the unrooted argument follow in
Sections~\ref{sec:helper}--\ref{sec:global-terminal}.
Appendix~\ref{sec:finite} contains the degree-seven lemma.

\clearpage
\section{Notation and external tools}\label{sec:preliminaries}

\subsection{Minors and densities}
Write $v(G)=|V(G)|$, $e(G)=|E(G)|$ and $\rho(G)=e(G)-4v(G)$.
A minor model consists of pairwise disjoint, nonempty connected
vertex sets, called \emph{bags}, with an edge between bags whenever
the target requires one. In a graph with prescribed roots $X$, a
rooted model retains each root in its corresponding bag, and no bag
contains two roots. Extra edges are allowed.
We allow an empty root set and interpret its rooted clique $K_0$ as
the empty model.

For a nonempty set $Y\subseteq V(G)$, let $N_G(Y)$ be its external
neighbourhood, and put
\[
 \rho_G(Y)=e(G[Y])+e_G(Y,N_G(Y))-4|Y|.
\]
We call $Y$ a \emph{fragment}, its neighbourhood its \emph{boundary},
and $V(G)\setminus(Y\cup N_G(Y))$ its \emph{opposite side}.
A fragment is \emph{proper} if its opposite side is nonempty.
In a rooted graph, fragments are required to avoid the roots. Set
\[
 \rho_4(G,X)=e(G)-e(G[X])-4|V(G)\setminus X|.
\]
Thus $\rho_4(G[Y\cup N_G(Y)],N_G(Y))=\rho_G(Y)$.

A $k$-rooted graph $(G,X)$ is \emph{$4$-light} if
$\rho_G(Y)\leq0$ for every root-free fragment with fewer than $k$
boundary vertices. In particular, for five roots the tested boundary
orders are at most four. The graph is \emph{internally $k$-connected}
if every nonempty root-free set has at least $k$ external neighbours.

Two disjoint, nonadjacent fragments form a \emph{bifragment}; it is
\emph{dense} if both fragments have positive density. Here nonadjacent
means that no edge joins the two fragments. An unrooted graph is
\emph{$4$-bilight} if it has no dense bifragment with both boundary
orders at most four. Every five-connected graph is $4$-bilight. Our main
extremal theorem is as follows.

\begin{theorem}\label{thm:bilight}
Every $4$-bilight graph $G$ with $v(G)\geq3$ and
$e(G)\geq4v(G)-2$ contains a $K_7^{-}$ minor.
\end{theorem}

A separation $(A,B)$ has $A\cup B=V(G)$ and no edge between
$A\setminus B$ and $B\setminus A$; its order is $|A\cap B|$.
It is proper if both open sides are nonempty. For a root separation
we also require $X\subseteq A$. Its right-hand side is $G[B]$
rooted at $A\cap B$. When a linkage carries these boundary roots
back to $X$, its paths are disjoint and their interiors avoid $B$.
Roots already on the boundary use trivial paths. Appending the paths
to their respective bags preserves disjointness and all contacts.

We repeatedly use the following lifting convention. Every vertex of
a minor is assigned its fixed connected preimage in the original
graph. Preimages of different vertices are disjoint. Replacing a bag
by the union of its preimages lifts connectivity, contacts and all
prescribed roots. Our minimal-counterexample arguments minimise
$(v(G),e(G))$ lexicographically; each reduction explicitly preserves
the relevant density and separation conditions.

\subsection{Rooted density theorems}
Let $\mathcal S_{5,j}$ be the class of graphs on five vertices with
at most $j$ edges, and let $\mathcal S_{5,4}^{-}$ omit
$K_3\mathbin{\dot\cup}K_2$ from $\mathcal S_{5,4}$.
The following form of Dvo\v r\'ak's theorem~\cite[Theorem~4]{Dvorak2026}
is reproduced in~\cite[Theorem~2.6]{DNR2026}.

\begin{theorem}\label{thm:rooted-targets}
If $(G,X)$ is a $4$-light five-rooted graph and $t=\rho_4(G,X)$,
then it has a rooted model of every labelled graph in $\mathcal T_t$,
where
\[
\begin{array}{c|cccccccc}
t&\leq0&1&2&3&4&5&6&\geq7\\ \hline
\mathcal T_t&\mathcal S_{5,0}&\mathcal S_{5,1}&\mathcal S_{5,3}&
\mathcal S_{5,4}^{-}&\mathcal S_{5,6}&\mathcal S_{5,8}&
\mathcal S_{5,9}&\mathcal S_{5,10}
\end{array}
\]
The labelling prescribes the root for each target vertex.
\end{theorem}

\begin{corollary}[{\cite[Lemma~5.6]{DNR2026}}]\label{thm:rooted-clique}
Let $(G,X)$ be $4$-light with five roots, and put $m=10-e(G[X])$.
If $\rho_4(G,X)\geq\lceil(m+3)/2\rceil$, then $G$ has a rooted
$K_5$ minor.
\end{corollary}
\begin{proof}
The target class at this threshold, and every higher threshold,
contains every graph with $m$ edges. Apply
Theorem~\ref{thm:rooted-targets} to the missing edges of $G[X]$,
then retain the existing root edges.
\end{proof}

A \emph{dart} has four roots, two root edges forming a three-vertex
path, and one nonroot adjacent to all four roots. The fourth root
is isolated in the root-induced graph. Dvo\v r\'ak's small-root
theorem~\cite[Corollary~13]{Dvorak2026}, together with
\cite[Lemma~3.1]{DNR2026}, gives the following fact: a $4$-light
$k$-rooted graph of positive rooted density contains a rooted $K_k$
if $k\leq3$, and a rooted dart if $k=4$.

\subsection{Reducible fragments and extraction}
A fragment $Y$ with boundary $S$ of order at most four is
\emph{reducible} if $\rho_G(Y)\leq0$ and $G[Y\cup S]$, rooted
at $S$, contains either a rooted clique on $S$, or a rooted dart
when $|S|=4$ and $|Y|\geq2$. For an unrooted ambient graph we
also require $|V(G)\setminus Y|\geq3$.

A \emph{reducent} replaces $Y$ by such a clique or dart, retaining
the graph outside $Y$. The replacement is a minor: the rooted model
provides fixed preimages for the boundary vertices and, for a dart,
its one additional vertex. In both cases the order strictly decreases.

\begin{lemma}[{\cite[Lemma~3.3]{DNR2026}}]\label{lem:reduction}
Replacing an inclusionwise-maximal reducible fragment preserves
$4$-bilightness and does not decrease $\rho$ in an unrooted
$4$-bilight graph. In a $4$-light five-rooted graph, it preserves
$4$-lightness and does not decrease $\rho_4$. All original roots
are retained in distinct prescribed bags.
\end{lemma}

For the extraction step, two weaker patterns from~\cite{DNR2026}
are useful. Each has five independent roots and two nonroots.
In the first, the nonroots are adjacent and each misses at most
one root, with distinct missed roots if both omissions occur.
In the second, the nonroots are nonadjacent and both meet every
root; this is $K_{2,5}$ rooted at its five-vertex part.

\begin{lemma}[{\cite[Corollary~4.2]{DNR2026}}]\label{lem:extraction}
Let $(G,X)$ be $4$-light, with $|X|\leq4$ and $\rho_4(G,X)\leq0$.
Suppose a root five-separation $(A,B)$ has a right-hand side containing
one of the two weaker patterns. Then $G$ has a reducible fragment
$Y$ containing $B\setminus A$.
\end{lemma}

The mechanism, also used for five roots below, is worth specifying.
Restrict either pattern to a proper subset $Z$ of its roots by deleting
the other root vertices. The remaining graph contains a rooted clique
if $|Z|\leq3$, or a dart if $|Z|=4$~\cite[Observation~4.1]{DNR2026}.
If five disjoint paths do not link the original roots to the boundary,
Menger's theorem gives an isolating separator of order at most four,
linked to a proper subset of the boundary. The clique or dart lifts
along this linkage. Its interior has at least two vertices, because
the two helper bags remain there. Lightness gives nonpositive density
when the isolator is smaller than the root set. When its order equals
the root count, choose the isolator $(X,V(G))$; its density is
$\rho_4(G,X)\leq0$ by assumption.
This gives the reducible fragment. When transferring it to an unrooted
ambient graph, we separately check that at least three vertices remain
outside it.

We also use the root-forgetting rule~\cite[Observation~5.2]{DNR2026}:
if $|X|\leq4$, $(G,X)$ is $4$-light and $\rho_4(G,X)\leq0$, then
the same is true after replacing $X$ by any subset $X'$. To see the
density bound, put $M=X\setminus X'$. For a set $Y$ tested in the
new rooted graph, or its entire nonroot set, put $Y'=Y\setminus M$.
The old hypotheses give $\rho_G(Y')\leq0$, interpreting the empty
set as having density zero. At most $4|M\cap Y|$ additional incident
edges are counted when $M\cap Y$ is restored: each such vertex has
possible neighbours, outside $Y'$, only in $M\cap Y$ or $N_G(Y)$,
and $|M|+|N_G(Y)|\leq4$. Hence
\[
 \rho_G(Y)\leq\rho_G(Y')+4|M\cap Y|-4|M\cap Y|\leq0.
\]

\section{Local rooted constructions}

We first record three small rooted constructions and a local contraction
lemma. They will also be used in the unrooted argument.

The rooted-star lemma is related to~\cite[Lemma~4.10]{DNR2026}, which
gives three spokes at one of five roots when $|V(H)|+|X|\le8$.
Here independence of $X$ allows all $k-1$ spokes and the bound
$|V(H)|+|X|\le2k-1$; the degree assumptions are stated below.

\begin{lemma}[Rooted star]\label{lem:star}
Let $Z\subseteq V(H)$ have order $k$, and let $X\subsetneq Z$ be
independent. Suppose that
\[
 |V(H)|+|X|\le 2k-1,\qquad
 d_H(v)\ge k-1\quad(v\notin X),\qquad
 d_H(x)\ge k-|X|\quad(x\in X).
\]
Then, for some $z\in Z\setminus X$, the set
$(V(H)\setminus Z)\cup\{z\}$ is connected and adjacent to every vertex
of $Z\setminus\{z\}$.
\end{lemma}

\begin{proof}
Put $T=V(H)\setminus Z$. A component $D$ of $H[T]$ misses at most
$|D|$ vertices of $Z$: any vertex of $D$ has at least $k-|D|$
neighbours in $Z$. Consequently at most $|T|$ roots miss some component
of $H[T]$. Since $|T|+|X|\le k-1$, there is a vertex
$z\in Z\setminus X$ adjacent to every component. Thus $T\cup\{z\}$
is connected, including when $T=\varnothing$.

A vertex $r\in Z\setminus\{z\}$ which has no neighbour in $T$ must
be adjacent to $z$. For $r\notin X$ this follows from $d_H(r)\ge k-1$;
for $r\in X$, independence of $X$ and $d_H(r)\ge k-|X|$ force
adjacency to every vertex of $Z\setminus X$.
\end{proof}

\begin{lemma}[Rooted triangle]\label{lem:triangle}
Let $(R,S)$ be internally three-connected, where $|S|=3$.
Suppose that $R-S$ is nonempty and connected, each root has a neighbour
outside $S$, and every nonroot has degree at least four. Then $R$
contains an $S$-rooted triangle minor.
\end{lemma}

\begin{proof}
Three distinct vertices $a,b,c$ of a two-connected graph have a rooted
triangle minor. To see this, take a cycle through $a,b$ and, if necessary,
a two-fan from $c$ to that cycle. Either their union contains a cycle
through all three vertices, or it contains a theta whose three paths
contain $a,b,c$ separately. In the latter case, assign one end of the
theta to the $a$-path and the other to the $b$-path; the interior of the
third path is the $c$-bag. These three bags are connected and pairwise
adjacent.

Consider the block-cut tree of $R$, representing a root which is a
cutvertex by its cutvertex node, and any other root by its block node.
If the median of the three root nodes is a block, the roots have three
distinct attachments to that block. The block is two-connected, so the
preceding observation supplies a triangle rooted at the attachments.
Appending the three disjoint exterior root paths gives the result.

If the median is a cutvertex $c$, each component of $R-c$ contains at
most one root. The nonroots in any such component would have boundary
contained in $c$ and that possible root, contrary to internal
three-connectivity. Hence $c$ is the only nonroot, but then its degree
is at most three. This is also impossible.
\end{proof}

We use the following terminal construction from the proof of
\cite[Lemma~3.1]{DNR2026}. Assume that four independent roots are the
leaves of a tree $T$. Its two degree-three branching vertices $v_1,v_2$
may coincide at a degree-four vertex. Choose $T$ to minimise the length
of the path between these vertices. If, after deleting $v_1,v_2$, every
partition of the roots into two pairs has a path joining the pairs,
then the graph has a rooted dart. This part of the cited proof uses
only the stated tree and path conditions, not a density hypothesis.

\begin{lemma}[Rooted dart]\label{lem:dart}
An internally four-connected graph with four prescribed roots, at least
two nonroots, and nonroot minimum degree at least five contains a rooted
dart.
\end{lemma}

\begin{proof}
Write $(F,Z)$ for the rooted graph and induct on $|V(F)|$. Delete the
root--root edges; this preserves all nonroot degrees and boundaries.
Every partition of $Z$ into two pairs has two disjoint paths between
the pairs. Otherwise Menger's theorem gives a separation of order at
most one, with one pair on each side. The nonroots on either open side
then have boundary contained in that pair and the separator, of order
at most three. Internal four-connectivity would put all nonroots in
the separator, contradicting their number.

Each component of $F-Z$ sees all four roots. It therefore supplies a
tree whose leaves are exactly $Z$. Choose such a tree with a shortest
branching path, with ends $v_1,v_2$. Unless $\{v_1,v_2\}$ separates
two roots from the other two, the terminal construction above gives a
dart. The case $v_1=v_2$ is covered, since a one-vertex separation of
this kind was already excluded.

In the remaining case let $(B_1,B_2)$ be the separation, with
$B_1\cap B_2=\{v_1,v_2\}$ and two roots in each open side. Put
\[
 Z_i=(Z\cap B_i)\cup\{v_1,v_2\},\qquad W_i=B_i\setminus Z_i.
\]
If both $W_i$ are empty, the degree bound makes $v_1,v_2$ adjacent to
each other and to every root. Use $v_1$ as the dart's helper and absorb
$v_2$ into one root bag: the resulting root star contains the required
three-vertex path.

Otherwise choose a nonempty $W_i$. It has at least two vertices, since
a single vertex would have at most the four neighbours in $Z_i$.
Every subset of $W_i$ has the same neighbours in $F[B_i]$ as in $F$.
Thus $(F[B_i],Z_i)$ is internally four-connected and retains the
nonroot degree bound. It omits the two opposite roots, so induction
provides a $Z_i$-rooted dart.

Take the two disjoint paths between the two original root pairs
established above. They cross the separator at different vertices.
Their segments after their last visits to $\{v_1,v_2\}$ lie outside
$B_i$, apart from their initial vertices, and end at the two opposite
roots. Append these segments to the bags at $v_1,v_2$. They are disjoint
from the other bags and from the helper, and restore the four original
roots without losing any contact.
\end{proof}

Call a graph \emph{quasi five-connected} if it is four-connected and has
no four-separation with at least two vertices on each open side. The
following local version of the crossing argument in
\cite[Lemma~1]{Kou2025} does not assume contraction-criticality.

\begin{lemma}\label{lem:quasi}
Let $Q$ be five-connected, and let $A$ have boundary $S=N_Q(A)$ of
order five. Suppose that $|A|\ge3$ and
$|V(Q)\setminus(A\cup S)|\ge2$. If $x\in S$ has a unique neighbour
$a$ in $A$, then $Q/xa$ is quasi five-connected.
\end{lemma}

\begin{proof}
An edge contraction preserves four-connectivity. A forbidden
four-separation of $Q/xa$ would lift to a five-cut $T$ of $Q$ containing
$x,a$, with open sides $B,D$ of order at least two. Write
$C=V(Q)\setminus(A\cup S)$ and
\[
 p=|A\cap T|,\quad r=|C\cap T|,\quad s=|S\cap T|,\quad
 u=|S\cap B|,\quad w=|S\cap D|.
\]
Then $p+r+s=u+w+s=5$ and $s\ge1$.

If $A\cap B\ne\varnothing$, its boundary is contained in
$(A\cap T)\cup(S\cap T)\cup(S\cap B)$ and omits $x$, whose unique
neighbour in $A$ lies in $T$. Five-connectivity gives $p+s+u\ge6$,
so $p>w$ and $u>r$. The opposite corner $C\cap D$ has boundary of
order at most $r+s+w<5$, and is therefore empty. Symmetrically,
$A\cap D\ne\varnothing$ implies $p>u$, $w>r$, and
$C\cap B=\varnothing$.

If both $A$-corners are nonempty, then $C\subseteq T$, so $r\ge2$,
whereas $u+w\ge2r+2\ge6>5-s$. If neither is nonempty, then
$p=|A|\ge3$ and $r\le1$. Some $C$-corner is nonempty, say $C\cap B$;
its boundary gives $u\ge p$. Hence $w\le1$ and $C\cap D$ is empty,
contrary to $|D|\ge2$.

Finally suppose only $A\cap B$ is nonempty. Then $D=S\cap D$, so
$w\ge2$. Since $p>w$, we have $r\le1$. But $C\cap B$ has boundary
of order at most $r+s+u=5+r-w\le4$, so it too is empty. This leaves
$|C|=r\le1$, the final contradiction.
\end{proof}

\section{Two rooted helpers}

For five prescribed roots $X$, say that a model has the
\emph{helper property} if it consists of five prescribed-root bags and
two further nonempty connected bags, all pairwise disjoint, and at least
ten of the eleven possible helper--root and helper--helper contacts
are present. The two helpers contain no root; no root--root contact is
required.

\begin{theorem}\label{thm:helper}
Every $4$-light five-rooted graph $(G,X)$ with $\rho_4(G,X)\ge2$ has
the helper property.
\end{theorem}

The proof first restricts a minimum counterexample. It then uses a
weighted neighbourhood to find an edge with few common neighbours.
An auxiliary clique allows Mader's atom theorem to be applied without
requiring intersections of positive-density fragments to remain positive.

\subsection{Minimum-counterexample reductions}

Suppose the theorem is false, and choose a counterexample minimising
$(|V(G)|,|E(G)|)$ lexicographically. Throughout this section, contractions
identify at most one original root in each preimage. A rooted model in
a minor lifts by replacing each of its vertices with the fixed connected
preimage; in particular, all five prescribed roots remain distinct.

We use the following consequence of the isolator argument in
\cite[Observation~4.1 and proof of Corollary~4.2]{DNR2026}. Suppose a
root-free region has a helper-property model on a five-vertex boundary,
or two helpers full to a larger boundary. Either five disjoint exterior
paths attach the original roots to five boundary bags, or a separator of
order $k\le4$ isolates the region and has a linkage to $k$ of its
boundary bags. A helper model restricted to those
$k$ bags contains a rooted $K_k$ for $k\le3$, and a rooted dart for
$k=4$. Indeed, retain a helper adjacent to all four roots and absorb the
other into a root bag to obtain two root contacts. For three roots,
absorb the helpers into two root bags; smaller cases follow by restriction.
The same conclusion holds when both helpers are full to a larger
boundary. If the isolated interior has at least two vertices and
nonpositive density, it is a reducible fragment. The linkage is outside
the original region, except for its ends, so it does not use a helper
vertex. This is the only form of the isolator argument needed below.

\begin{lemma}\label{rh:basic}
The minimum counterexample has the following properties:
\begin{enumerate}
\item $X$ is independent, and there is no reducible root-free fragment
with boundary of order at most four.
\item Every proper root-free set $Y$ with $|N_G(Y)|=5$ satisfies
$\rho_G(Y)\le1$.
\item $G-X$ is connected, every nonroot has at most three root neighbours,
and every root has degree at least two.
\item $\rho_4(G,X)=2$.
\end{enumerate}
\end{lemma}

\begin{proof}
Root--root edges contribute neither to the density nor to a required
helper contact, and cannot occur inside a bag containing just one root.
Thus minimality makes $X$ independent. A maximal reducible fragment
would give, by \cite[Lemma~3.3]{DNR2026}, a smaller $4$-light rooted
minor of nondecreasing density. Its helper model would lift, proving
(1).

For (2), put $S=N_G(Y)$. The induced rooted graph $(G[Y\cup S],S)$
is $4$-light, since its root-free sets retain their incident edges and
neighbours. If its density is at least two, it is a smaller graph with
a helper model by minimality. Five exterior root paths would extend
that model to $X$. Otherwise the isolator argument gives a reducible
fragment: its interior contains the two helpers and has nonpositive
density by $4$-lightness. Both possibilities contradict (1).

The densities of the components of $G-X$ sum to $\rho_4(G,X)$. Each
positive-density component sees all five roots, by $4$-lightness. Two
such components give two full helpers, with only their mutual contact
possibly missing. Hence there is exactly one positive component and
its density is at least two. Any other component would make this a
proper five-boundary set, contrary to (2). Thus $G-X$ is connected and
sees every root.

Suppose a nonroot $z$ has $d\ge4$ root neighbours. For the components
$Q_i$ of $G-(X\cup\{z\})$,
\[
 \sum_i\rho_G(Q_i)=\rho_4(G,X)+4-d.
\]
Each $Q_i$ sees $z$, and each positive one sees at least four roots.
If $d=5$, there is a positive component; it and $\{z\}$ are adjacent
helpers with at least ten contacts. If $d=4$, let $x$ be the missing
root. A component seeing all roots already gives the required model.
Otherwise each positive component has exactly four root neighbours,
and has density at most one by (2). Their sum is at least two, so there
are at least two positive components. Choose a component $R$ seeing
$x$, and a positive component $Q\ne R$. The helpers $\{z\}\cup R$
and $Q$ are adjacent and have respectively five and at least four root
contacts. This again contradicts the choice of $G$.

If a root $x$ had just one neighbour $v$, the side $G-x$, rooted at
$(X\setminus\{x\})\cup\{v\}$, would have density
$\rho_4(G,X)+4-|N_G(v)\cap X|\ge3$. By $4$-lightness, all five
proposed boundary vertices must be actual neighbours of its interior.
It is therefore a proper five-boundary side, contradicting (2). This
proves (3).

Finally suppose $\rho_4(G,X)\ge3$ and delete an edge with a nonroot
end. The result has density at least two but, by minimality, is not
$4$-light. A positive small-boundary set $Y$ becomes permissible only
because the deleted edge joins $Y$ to a vertex outside its closed
side. Restoring that edge gives a five-boundary set of density at least
two. By (2) its closed side is all of $G$, and its boundary is $X$.
The restored root endpoint then has only one neighbour, contrary to
(3). Therefore the density is exactly two.
\end{proof}

\begin{lemma}\label{rh:internal}
Every nonroot has degree at least five, and $(G,X)$ is internally
five-connected.
\end{lemma}

\begin{proof}
Deleting a nonroot $v$ of degree at most four leaves density
$2+4-d_G(v)\ge2$. If $G-v$ has a positive root-free set $Y$ with at
most four neighbours, restoring $v$ must add a neighbour to $Y$;
otherwise $Y$ already violates $4$-lightness in $G$. Its restored
density is at least two and its boundary has order five. This is a
proper side, since its boundary contains the nonroot $v$ and could not
also contain all five roots. Lemma~\ref{rh:basic}(2) excludes it.
Thus deletion stays in the induction class, a contradiction.

If internal five-connectivity fails, choose a connected root-free set
$Y$ of minimum boundary order $k\le4$, and put $S=N_G(Y)$. The rooted
shore $(G[Y\cup S],S)$ is internally $k$-connected, retains all nonroot
degrees, and has nonpositive density. For $k=0,1,2$, it contains a rooted
$K_k$ (using a path through $Y$ when $k=2$). For $k=3$, apply
Lemma~\ref{lem:triangle}. For $k=4$, the degree bound implies $|Y|\ge2$,
and Lemma~\ref{lem:dart} applies. In every case $Y$ is reducible,
contrary to Lemma~\ref{rh:basic}(1).
\end{proof}

Complete $X$ to a clique, obtaining $Q_0=G+K_X$. It is five-connected:
after deleting at most four vertices the surviving roots lie in one
clique, and another component would contradict internal five-connectivity.
For an original edge $e$ having a nonroot end, let $t(e)$ be its number
of common neighbours in $Q_0$. The rooted density after contraction is
\begin{equation}\label{rh:contract-density}
 \rho_4(G/e,X_e)=2+3-t(e),
\end{equation}
where $X_e$ denotes the unchanged or naturally identified roots. The
completion counts exactly the edges which become root--root edges when
one endpoint is a root.

\begin{lemma}\label{rh:blocker}
If $e=uv$ has a nonroot end and $t(e)\le3$, there is a proper root-free
set $Y$ such that
\[
 |N_G(Y)|=5,\qquad \rho_G(Y)=1,\qquad \{u,v\}\subseteq N_G(Y),
\]
and no vertex of $Y$ is adjacent to both $u$ and $v$.
\end{lemma}

\begin{proof}
By \eqref{rh:contract-density} and minimality, $G/e$ is not $4$-light.
Let $Y$ be a positive root-free set witnessing this. The identified
vertex cannot lie outside its closed side, because then $Y$ is a
root-free set of boundary at most four in $G$. Nor can it lie inside
$Y$: its preimage would have the same small boundary (and would be
root-free; a contracted root cannot be interior). It therefore lies
on the boundary. Splitting it gives exactly five neighbours in $G$.
The side is proper, since its boundary contains a nonroot and hence
cannot contain all five original roots. If $c$ vertices of $Y$ see
both endpoints, then
$\rho_G(Y)=\rho_{G/e}(Y)+c$. Lemma~\ref{rh:basic}(2) forces both
densities to equal one and $c=0$.
\end{proof}

\begin{lemma}\label{rh:five-edge}
There are no adjacent degree-five nonroots with exactly three common
neighbours.
\end{lemma}

\begin{proof}
Suppose $p,q$ are such vertices. Write their common neighbours as the
three-set $A$, and their respective exclusive neighbours as $u,v$.
Lemma~\ref{rh:blocker} applied to $pq$ gives a density-one set $T$
with boundary $\{p,q\}\cup D$, where $|D|=3$, avoiding $A$.
Since both $p,q$ have neighbours in $T$, their neighbourhoods force
$u,v\in T$; in particular these are nonroots.

Delete $p,q$ from this shore and reroot at $\{u,v\}\cup D$:
\[
 R=G[T\cup D],\qquad Z=\{u,v\}\cup D.
\]
The graph $(R,Z)$ is $4$-light. Indeed, every subset of its nonroots
$T\setminus\{u,v\}$ has exactly the same neighbours and incident
edges as in $G$, since only $u,v$ in $T$ see $p,q$.
Put $a_0=e_G(\{u,v\},D)$ and $\delta=1$ if $uv\in E(G)$, and
$\delta=0$ otherwise. Then
\[
 \rho_4(R,Z)
 =(4|T|+1-2-\delta-a_0)-4(|T|-2)
 =7-\delta-a_0.
\]
This is exactly the number $m\in\{0,\ldots,7\}$ of missing edges
from $uv$ together with the six edges between $\{u,v\}$ and $D$.
The rooted universality theorem \cite[Theorem~2.6]{DNR2026} supplies
all these missing edges simultaneously: its table contains every
$m$-edge graph in this range (the exceptional graph at density three
has four edges). Retain the root edges already present. We obtain
five rooted bags in $R$, with the bags at $u,v$ adjacent to one another
and to all three bags at $D$.

Append $p$ to the $u$-bag and $q$ to the $v$-bag. These are adjacent
root-free helpers, full to the boundary $A\cup D$. For $A\cap D$
use the existing $D$-bags; for $A\setminus D$ use singletons. These
choices are disjoint because $T\cap A=\varnothing$.

The root-free set $U=T\cup\{p,q\}$ has boundary exactly $A\cup D$.
If five exterior paths link $X$ to this boundary, retain the five bags
at their distinct ends and append the paths. This gives the helper
property rooted at $X$. Otherwise the isolator argument gives a
separator of order at most four, a rooted clique or dart on it, and a
root-free interior containing $U$. That interior has at least two
vertices and nonpositive density, so is reducible. Both alternatives
are impossible.
\end{proof}

\subsection{Eliminating degree five}

\begin{lemma}\label{rh:quasi-light}
If $u,v$ are nonroots and $Q_0/uv$ is quasi five-connected, then
$(G/uv,X)$ is $4$-light.
\end{lemma}

\begin{proof}
Root completion does not change any root-free boundary or incident
density. Every root has degree at least five in $Q_0/uv$: its at least
two nonroot neighbours merge to at least one, and it retains the four
other roots as neighbours. A root-free set with boundary at most four
has nonempty opposite side, because there are five roots. Its boundary
therefore has order four, and quasi five-connectivity makes one open
side a singleton. A singleton root-free side has degree four and
density zero. A singleton opposite side would be a root of degree at
most four, which is impossible.
\end{proof}

\begin{lemma}\label{rh:degree-six}
Every nonroot of $G$ has degree at least six.
\end{lemma}

\begin{proof}
Suppose a nonroot $v$ has degree five. First assume $Q_0[N_G(v)]$ is
not complete, and choose a nonedge $ab$ in it. Its ends are not both
roots. The graph
\[
 H=G-v+ab
\]
has rooted density two and is a proper rooted minor: contract $va$
and delete the unwanted new edges. If $H$ is $4$-light, its helper
model lifts, contrary to minimality.

Otherwise take a positive root-free set $Y$ of $H$ with at most four
neighbours. Put $k=e_G(v,Y)$ and let $c$ indicate whether the added
edge $ab$ is incident with $Y$. Then
\[
 N_G(Y)\subseteq N_H(Y)\cup\{v\},\qquad
 \rho_G(Y)=\rho_H(Y)+k-c.
\]
Internal five-connectivity gives $|N_G(Y)|=5$ and $k\ge1$. This side
is proper, since its boundary contains the nonroot $v$. Consequently
\[
 1\ge\rho_G(Y)=\rho_H(Y)+k-c\ge1+k-c,
\]
forcing $\rho_G(Y)=\rho_H(Y)=k=c=1$. Exactly one endpoint of $ab$
lies in $Y$; call it $a$. Then $a$ is a nonroot and the unique neighbour
of $v$ in $Y$. The other endpoint $b$ is among the four vertices of
$N_G(Y)\setminus\{v\}$, since the added edge cannot introduce a fifth
neighbour in $H$.

The set $Y$ is not a singleton, which would make $ab$ an edge of $G$.
If $|Y|=2$, its density one and the nonroot degree bound force adjacent
degree-five vertices with three common neighbours, contrary to
Lemma~\ref{rh:five-edge}. Hence $|Y|\ge3$. Its opposite side in
$Q_0$ has at least two vertices: a singleton would be an original root
$z$, with boundary $(X\setminus\{z\})\cup\{v\}$, and independence
of $X$ would give $d_G(z)\le1$.

Lemma~\ref{lem:quasi} now applies to the unique edge $va$ into $Y$.
Thus $Q_0/va$ is quasi five-connected, and Lemma~\ref{rh:quasi-light}
makes $G/va$ $4$-light. Since $v$ has degree five and $b$ misses $a$,
the edge $va$ has at most three common neighbours. Its contraction
retains rooted density at least two and decreases order, a contradiction.

It remains that $Q_0[N_G(v)]$ is a clique. Internal five-connectivity
and Menger's theorem give five disjoint paths from $X$ to $N_G(v)$,
with interiors outside $N_G[v]$. Roots already in $N_G(v)$ use trivial
paths. Choose $z\in N_G(v)\setminus X$, which exists because $v$ has
at most three root neighbours. Use $\{v\},\{z\}$ as helpers, and the
five linkage paths as root bags, removing $z$ from its own path. That
trimmed bag is nonempty because $z$ is not an original root. The helper
$z$ meets all five root bags, $v$ meets the other four, and $vz$ joins
the helpers. Every used clique edge has nonroot end $z$, so is an actual
edge of $G$. This gives the ten required contacts.
\end{proof}

\subsection{Weighted neighbourhoods and an atom}

Let $E_0$ be the original edges with a nonroot end and at most three
common neighbours in $Q_0$. Each edge of $E_0$ has a five-boundary
set as in Lemma~\ref{rh:blocker}.

\begin{lemma}\label{rh:weighted}
If $v$ is a nonroot and $d_G(v)+|N_G(v)\cap X|\le9$, then some edge
of $E_0$ is incident with $v$.
\end{lemma}

\begin{proof}
Suppose instead that $t(vu)\ge4$ for every neighbour $u$ of $v$.
Put $H=G[N_G(v)]$, $X_0=N_G(v)\cap X$, and $r=|X_0|$. Vertices of
$H-X_0$ have degree at least four in $H$; vertices of $X_0$ have degree
at least $5-r$, since root completion adds exactly $r-1$ common
neighbours to such an edge. Also $X_0$ is independent and
$|V(H)|+r\le9$.

There are five disjoint paths from $X$ to $N_G(v)$ in $G-v$.
Otherwise a separator of order at most four would isolate, in $G$,
a root-free component containing $v$. Truncate the paths at their first
entries to $N_G(v)$. Their five ends form a set $Z$ containing $X_0$,
and their interiors avoid the neighbourhood. Lemma~\ref{lem:star}
gives a central bag
\[
 C=(N_G(v)\setminus Z)\cup\{z\},\qquad z\in Z\setminus X_0,
\]
adjacent to the four other ends. Use $C$ and $\{v\}$ as helpers,
keeping the four other linkage paths and trimming $z$ from its path.
The trimmed root bag is nonempty; all bags are disjoint. The helper
$C$ has five root contacts, $v$ has four, and they are adjacent. This
is the helper property, a contradiction.
\end{proof}

In particular, $E_0$ is nonempty. If $m=|V(G)\setminus X|$, the
nonroot degree bound and at most three root neighbours give $m\ge4$,
while
\begin{equation}\label{rh:total-weight}
 \sum_{v\notin X}\bigl(d_G(v)+|N_G(v)\cap X|\bigr)=8m+4.
\end{equation}
Some nonroot therefore has weight at most nine.

We use Mader's atom theorem in the form stated in
\cite[Theorem~5]{KS2017}. In a graph $Q$ of connectivity $k$, let
$\mathcal S$ be any family of vertex sets. An $\mathcal S$-fragment
is a nonempty proper union of components of $Q-T$, where $T$ is a
$k$-cut containing a member of $\mathcal S$. An $\mathcal S$-atom
is such a fragment of minimum order. For an atom $A$, put
$T_A=N_Q(A)$ and $\overline A=V(Q)\setminus(A\cup T_A)$. If some
$S\in\mathcal S$ and $k$-cut $T$ satisfy
\begin{equation}\label{rh:atom-hyp}
 S\subseteq T\setminus\overline A,\qquad A\cap T\ne\varnothing,
\end{equation}
then
\begin{equation}\label{rh:atom-conclusion}
 A\subseteq T,\qquad |A|\le\tfrac12|T\setminus T_A|.
\end{equation}
There is no assumption that every edge is noncontractible.

\begin{proof}[Completion of the proof of Theorem~\ref{thm:helper}]
Add a set $W$ of $|V(G)|+1$ vertices and make $X\cup W$ a clique,
with no edges from $W$ to $V(G)\setminus X$; call the resulting graph
$Q$. Internal five-connectivity shows that $Q$ is five-connected, and
$X$ separates $G-X$ from $W$, so its connectivity is exactly five.
This padding changes neither boundaries of original root-free sets nor
common-neighbour counts of original edges with a nonroot end.

Let $\mathcal S=\{V(e):e\in E_0\}$. By Lemma~\ref{rh:blocker},
$\mathcal S$-fragments exist in $Q$, with an example of order at most
$m=|V(G)\setminus X|$. Choose an $\mathcal S$-atom $A$, among all
such fragments, so $|A|\le m$. It avoids $X\cup W$: otherwise its
five-boundary would force it to contain the whole surviving clique,
of order at least $|W|>m$. Thus $A\subseteq V(G)\setminus X$, and
its boundary is the same in $G$ and $Q$.

Its closed side is proper in $G$. If not, its five-boundary would be
exactly $X$, and could not contain the ends of an edge in $E_0$.
Lemma~\ref{rh:basic}(2) therefore gives $\rho_G(A)\le1$. The
nonroot degree bound excludes $|A|=1,2$, since for these orders
\[
 \rho_G(A)\ge2|A|-\binom{|A|}{2}\ge2.
\]
Writing $S_A=N_G(A)$, double-counting incident edges gives
\[
 \sum_{v\in A}\bigl(d_G(v)+|N_G(v)\cap S_A|\bigr)
   =8|A|+2\rho_G(A)\le8|A|+2.
\]
As $|A|\ge3$, some vertex $v\in A$ has the displayed weight at most
eight. Its root weight is no larger, so Lemma~\ref{rh:weighted}
supplies an incident edge $f\in E_0$. Both ends of $f$ lie in
$A\cup T_A$.

Apply Lemma~\ref{rh:blocker} to $f$, and let $T$ be the boundary of
its resulting fragment. In $Q$, this is a five-cut containing both
ends of $f$ and meeting $A$ at $v$. Consequently
$V(f)\subseteq T\setminus\overline A$, so \eqref{rh:atom-hyp} holds.
Mader's theorem yields $|A|\le5/2$, contradicting $|A|\ge3$.

The auxiliary clique has only selected an atom: no model has been
constructed in the padded graph. The contradiction proves the theorem
in the original rooted graph.
\end{proof}

\section{Orienting the separations}
\label{sec:global-orientations}

Suppose that Theorem~\ref{thm:bilight} fails, and choose a
counterexample $G$ minimising $(v(G),e(G))$ lexicographically. Thus $G$
is $4$-bilight, has at least three vertices and $\rho(G)\geq-2$, and
has no $\Kminus$ minor. Simplicity gives $v(G)\geq9$. Throughout the
next two sections, $G$ denotes this minimum counterexample.

For a fragment $A$, write $R_A$ for $G[A\cup N_G(A)]$ rooted at
$N_G(A)$. We call a five-rooted side \emph{heavy} if its rooted
density is at least two, and \emph{low} if its density is at most one.
For a separation $(A,B)$, write $R_{A,B}=G[B]$, rooted at $A\cap B$.

\begin{lemma}\label{gl:reducible}
The graph $G$ has no reducible fragment.
\end{lemma}

\begin{proof}
Replace an inclusion-maximal reducible fragment by a clique or dart
rooted at its boundary. By Lemma~\ref{lem:reduction}, the resulting
minor is $4$-bilight and has density at least $\rho(G)$. The
definition of reducibility leaves at least three vertices outside the
fragment. The replacement strictly decreases the order: a clique
removes every interior vertex, and a dart retains one where there were
at least two. This contradicts minimality.
\end{proof}

The next lemma adapts the separation argument
of~\cite[Section~5]{DNR2026}. We give the details to account for
the density threshold of two and the final $\Kminus$ model.

\begin{lemma}\label{gl:orientations}
The following assertions hold.
\begin{enumerate}
\item If one side of a five-separation is heavy, the opposite side
is $4$-light.
\item For every five-separation with boundary $X$, exactly one side
is heavy. If $m=10-e(G[X])$, the heavy side has density at least $m+7$.
\item There is no bifragment $(S,T)$ with both boundaries of order
at most five, both densities positive, and every member with boundary
of order five heavy.
\end{enumerate}
\end{lemma}

\begin{proof}
We first establish an observation about small separations. Suppose
$(C,D)$ has order at most four, $R_{C,D}$ is $4$-light and has
nonpositive density, and
\[
 |(C\setminus D)\cup N_G(C\setminus D)|\geq3.
\]
There cannot be a five-separation $(A,B)$ with $B\subseteq D$ and
$R_{A,B}$ both $4$-light and heavy. To see this, remove from $C\cap D$
the vertices with no neighbour in $C\setminus D$, obtaining $C'$.
The rooted graph $H=R_{C',D}$ remains $4$-light and has nonpositive
density, by the root-forgetting observation
\cite[Observation~5.2]{DNR2026}: if a root set of order at most four
is reduced, each forgotten root contributes at most four incident
edges outside a given old fragment, offset by its density cost of four.
The roots of $H$ lie in $A$, so $(A\cap D,B)$ is a root separation
of $H$ with right side $R_{A,B}$. The helper model supplied there by
Theorem~\ref{thm:helper} contains one of the rooted graphs in
Lemma~\ref{lem:extraction}. That lemma gives a reducible fragment
of $H$, which is also reducible in $G$: at least $|C'|\geq3$
vertices remain outside it. This contradicts Lemma~\ref{gl:reducible}.

For (1), suppose that the side opposite a heavy side $R_{A,B}$ is
not $4$-light. It contains a positive fragment $L\subseteq A\setminus B$
with at most four neighbours. Bilightness implies that $R_{A,B}$ is
$4$-light, since a positive small-boundary fragment inside $B\setminus A$
would form a dense bifragment with $L$. Set
\[
 (C,D)=(L\cup N_G(L),\,V(G)\setminus L).
\]
The side $R_{C,D}$ is $4$-light and has nonpositive density: otherwise
its whole interior, or a positive small-boundary fragment inside it,
would again form a dense bifragment with $L$. Also
$|L\cup N_G(L)|\geq6$, since $L$ has positive density. Now
$B\subseteq D$, contradicting the preceding observation.

Let $a\geq b$ be the two densities at a five-separation with boundary
$X$. Then
\begin{equation}\label{gl:density-sum}
 a+b=\rho(G)+20-e(G[X])\geq m+8.
\end{equation}
In particular, $a\geq\lceil(m+8)/2\rceil$. If $b\geq2$, (1) makes
both sides $4$-light. The rooted clique bound
(Corollary~\ref{thm:rooted-clique}) applies to the side of density $a$;
Theorem~\ref{thm:helper} applies to the other. Join corresponding root
branch sets at their common root. The five resulting branch sets
are pairwise adjacent; the two helper branch sets remain in the second
open side and retain at least ten of their eleven possible adjacencies.
This gives a $\Kminus$ model. Thus
$b\leq1$, and \eqref{gl:density-sum} gives $a\geq m+7$, proving (2).

For (3), first consider two heavy five-fragments $S,T$ forming a
bifragment. By (1), both $R_S$ and $R_T$ are $4$-light. We claim that
five disjoint paths join their closed sides. Otherwise
Menger's theorem gives a separation $(C,D)$ of order at most four
containing $S\cup N_G(S)$ in $C$ and $T\cup N_G(T)$ in $D$.
Here $|C\setminus D|\geq3$, since $|S|\geq2$.

We verify that $R_{C,D}$ is $4$-light and has nonpositive density.
Discard from its root set the vertices with no neighbour in
$D\setminus C$, leaving a separation $(C,D')$. No vertex of $S$
can lie in $C\cap D'$, since all neighbours of $S$ lie in $C$.
Consequently $R_{C,D'}$ is a root side of order at most four in the
side opposite $R_S$. By (1) it has nonpositive density, equal to
that of $R_{C,D}$. Applying the same argument to every root side of
$R_{C,D}$ proves its $4$-lightness. The initial observation, applied
to $R_T$, gives a contradiction. The five paths therefore exist.

Trim these paths so that they run between $N_G(S)$ and $N_G(T)$,
with interiors outside both closed sides. The ends are distinct on
each five-element boundary; a shared boundary vertex uses a trivial
path. Apply the helper theorem in $R_S$ and the rooted clique bound
in $R_T$, using (2). Join corresponding root branch sets along these
paths. Their interiors are disjoint; assigning each to its corresponding
branch set preserves disjointness and connectivity. Again they give a
$\Kminus$ model.
This excludes two heavy five-fragments. A heavy five-fragment and a
positive fragment of boundary at most four contradict (1). Two
positive fragments of boundary at most four contradict bilightness.
\end{proof}

\begin{lemma}\label{gl:exact-density}
We have $e(G)=4v(G)-2$.
\end{lemma}

\begin{proof}
If $\rho(G)\geq-1$, delete an edge $uv$ to obtain $J$. Minimality
forces a dense bifragment $(S,T)$ in $J$, with both boundaries of
order at most four. Unless this already violates bilightness in $G$,
we may label the ends so that $u\in S$ and
$v\notin S\cup N_J(S)$. If $v\in T$, choose the labels also so that
$|N_G(S)|\geq|N_G(T)|$. Then $|N_G(S)|\leq5$ and
$\rho_G(S)=\rho_J(S)+1\geq2$, so $|S|\geq2$.

Put $S'=S\setminus\{u\}$. We have
\[
 N_G(S')\subseteq(N_G(S)\setminus\{v\})\cup\{u\},\qquad
 \rho_G(S')=\rho_G(S)+4-e_G(\{u\},N_G(S)).
\]
The pair $(S',T)$ is nonadjacent in $G$. If $|N_G(S)|\leq4$, both
members have boundary at most four and positive density, a
contradiction. Otherwise Lemma~\ref{gl:orientations}(2) gives
$\rho_G(S)\geq7$, hence $\rho_G(S')\geq6$. The density of $T$ is
positive, and at least two if its boundary grows to order five.
Lemma~\ref{gl:orientations}(3) again gives a contradiction.
\end{proof}

\section{Degree bounds and five-vertex cuts}
\label{sec:global-cuts}

\begin{lemma}\label{gl:five-connected}
The graph $G$ has minimum degree at least five and is five-connected.
\end{lemma}

\begin{proof}
If $d_G(v)\leq4$, the minor $J=G-v$ has density at least $-2$
and at least eight vertices. Any dense bifragment in $J$ with
boundaries of order at most four remains a bifragment in $G$. Its
members remain positive; if a boundary grows to order five, the
restored vertex adds at least one incident edge, so the density is
at least two. Lemma~\ref{gl:orientations}(3) excludes such a pair.
Thus $J$ is $4$-bilight, contrary to minimality.

Suppose that $G$ has connectivity $k\leq4$, and let $S$ be a minimum
cut, possibly empty. Each component of $G-S$ has neighbourhood $S$.
By bilightness some component $A$ has nonpositive density. The side
$R_A$ is internally $k$-connected, and its nonroots retain degree at
least five. For $k\leq2$, its connected interior gives a rooted
clique. For $k=3$, apply Lemma~\ref{lem:triangle}; for $k=4$,
apply Lemma~\ref{lem:dart}. In the last case $|A|\geq2$, since the
minimum degree is five. Another component together with $S$ has at
least six vertices, again by the degree bound. Hence $A$ is reducible,
contrary to Lemma~\ref{gl:reducible}.
\end{proof}

Write $t(xy)=|N_G(x)\cap N_G(y)|$ for an edge $xy$. Contracting this edge
changes the global density by
\begin{equation}\label{gl:contraction-density}
 \rho(G/xy)=\rho(G)+3-t(xy).
\end{equation}
A quasi-five-connected graph is $4$-bilight: any member of a dense
bifragment has boundary four and at least two vertices, and the
other member supplies at least two vertices on the opposite side of
that cut. Thus a quasi-five-connected contraction with $t(xy)\leq3$
would contradict minimality.

We will repeatedly use the following observation. Contracting an edge
of a five-connected graph gives a four-connected graph, and every
four-cut contains the merged vertex. Consequently, in a dense
bifragment of the contraction with boundaries of order at most four,
the merged vertex belongs to both boundaries. Its open sides lift to
positive five-fragments of the original graph, with both contracted
ends in each boundary. Their densities do not decrease; they remain
disjoint and nonadjacent. Each has at least two vertices, because a
singleton of boundary four has density zero in the contraction.

\begin{lemma}\label{gl:degree-six}
The minimum degree of $G$ is at least six.
\end{lemma}

\begin{proof}
Suppose that $d_G(v)=5$. If $N_G(v)$ is a clique, choose a component
$C$ outside $N_G[v]$. Such a component exists because $v(G)\geq9$,
and five-connectivity gives $N_G(C)=N_G(v)$. The five singleton
neighbours, $\{v\}$ and $C$ form a $\Kminus$ model.

Otherwise choose nonadjacent $a,b\in N_G(v)$ and set $J=G-v+ab$.
This is a minor: contract $va$, then delete the new edges at $a$
other than $ab$. It has density $-2$, so a dense bifragment $(Y,Z)$
with boundaries of order at most four exists in $J$. For either
member, say $Y$, put $k=e_G(\{v\},Y)$ and let $c$ be one if the
added edge $ab$ has an end in $Y$, and zero otherwise. Then
\begin{equation}\label{gl:split-density}
 \rho_G(Y)=\rho_J(Y)+k-c.
\end{equation}
Since $c=1$ implies $k\geq1$, both sets remain positive in $G$;
they are still disjoint and nonadjacent, and their boundaries have
order at most five. By Lemma~\ref{gl:orientations}(3), we may choose
$Y$ with boundary five and $\rho_G(Y)=1$. Its boundary
grows on restoring $v$, so $k\geq1$. Equation~\eqref{gl:split-density}
forces $\rho_J(Y)=k=c=1$. Exchange $a,b$ if necessary: $a$ is the
unique neighbour of $v$ in $Y$, and $b$ lies in the boundary of $Y$
in both graphs. Hence $|Y|\geq2$. The opposite side of $Y$ in $G$
contains $Z$, which has at least two vertices.

If $|Y|\geq3$, Lemma~\ref{lem:quasi} makes $G/va$
quasi-five-connected. As $ab$ is absent and $d_G(v)=5$, we have
$t(va)\leq3$, a contradiction. We are left with $Y=\{a,c\}$.
The equation
\[
 1=d_G(a)+d_G(c)-e(G[\{a,c\}])-8
\]
forces $ac\in E(G)$ and $d_G(a)=d_G(c)=5$. Their common neighbours
form a triple $R$, and their complete neighbourhoods are
\begin{equation}\label{gl:first-pair}
 N_G(a)=\{c,v\}\cup R,\qquad
 N_G(c)=\{a,b\}\cup R,\qquad |R|=3.
\end{equation}
In particular $vb$ is an edge.

Contract $ac$, which has exactly three common neighbours. By
minimality and the preceding contraction observation, a dense
bifragment lifts to two positive five-fragments. Choose a member
$T$ of density one, using Lemma~\ref{gl:orientations}(3). Lifting
increases its density by $|T\cap R|$, so $T$ avoids $R$. Both $a,c$
are in its boundary; thus $v,b\in T$, and $v$ is the unique
neighbour of $a$ in $T$. Its opposite side has at least two vertices.
If $|T|\geq3$, Lemma~\ref{lem:quasi} applies to $av$; the missing
edge $cv$ gives $t(av)\leq3$, again a contradiction. Therefore
$T=\{v,b\}$, and the same degree calculation gives
\begin{equation}\label{gl:second-pair}
 N_G(v)=\{a,b\}\cup D,\qquad
 N_G(b)=\{c,v\}\cup D,\qquad |D|=3.
\end{equation}
The four degree-five vertices induce the cycle $a,c,b,v,a$, and
$R,D$ avoid the cycle.

If $|R\cap D|\geq2$, the cycle has at most four external neighbours.
There is a vertex outside the cycle and its boundary, because
$v(G)\geq9$, contradicting five-connectivity. If $|R\cap D|=1$,
put $S=R\cup D$ and $U=\{a,c,b,v\}$. Here $|S|=5$ and
$\rho_G(U)=0$. Its opposite side $W$ has
\[
 \rho_G(W)=18-e(G[S])\geq8.
\]
In particular $W$ is nonempty. Its closed side rooted at $S$ is
internally five-connected, so the rooted clique theorem applies.
The connected branch sets $\{a,v\}$ and $\{c,b\}$ are adjacent to each
other and to every vertex of $S$. Together with the rooted clique
they give a $K_7$ model.

Finally suppose that $R\cap D=\varnothing$. Every five-cut $L$
containing $a,v$ must also contain $c,b$. Indeed, if $c\notin L$,
all of $R$ lies in its component of $G-L$ or in $L$. Then $a$ has
no neighbour in any other component, and removing $a$ from $L$
leaves a cut of order four. The argument for $b$ is the same, using
the neighbours of $v$.

Now $t(av)=0$. The contraction $G/av$ satisfies the density bound, so
it has a dense bifragment. Both lifted five-boundaries contain all
four cycle vertices, by the preceding paragraph. Each open side
meets both $R$ and $D$, because its boundary contains $a,v$.
One of these disjoint sides, say $B$, contains exactly one vertex
$r$ of $R$. It cannot be a singleton, since its image has positive
density and boundary four. If $|B|=2$, then $B=\{r,d\}$ for some
$d\in D$, and its boundary consists of the four cycle vertices
and one further vertex $z$. The vertex $r$ is adjacent to neither
$v$ nor $b$, so its neighbours would all lie in $\{a,c,d,z\}$,
contrary to the minimum degree five.

Thus $|B|\geq3$, its opposite side has at least two vertices,
and $a$ has the unique neighbour $r$ in $B$. Lemma~\ref{lem:quasi}
makes $G/ar$ quasi-five-connected. Also $t(ar)\leq3$, since $a$
has degree five and its neighbour $v$ is not adjacent to $r$.
This final contraction contradicts minimality. Every contraction
above leaves at least eight vertices, preserving the order condition.
\end{proof}

\begin{lemma}\label{gl:low-edge}
Every proper low five-fragment $A$ contains an endpoint of an edge
$e$ with $t(e)\leq3$.
\end{lemma}

\begin{proof}
Put $S=N_G(A)$ and $a=|A|$. The minimum degree six gives
\[
 \rho_G(A)=\sum_{v\in A}d_G(v)-e(G[A])-4a
 \geq2a-\binom a2,
\]
so $a\geq5$. Also
\[
 \sum_{v\in A}\bigl(d_G(v)+|N_G(v)\cap S|\bigr)
 =8a+2\rho_G(A)\leq8a+2.
\]
Choose $v\in A$ whose summand is at most eight. Suppose that every
edge at $v$ has at least four common neighbours. Let $H$ be
$G[N_G(v)]$ with the edges within $X_0=N_G(v)\cap S$ deleted.
Vertices of $H-X_0$ have degree at least four, and vertices in
$X_0$ have degree at least $5-|X_0|$. Moreover
$v(H)+|X_0|\leq8$.

The side $G[A\cup S]$, rooted at $S$, is internally five-connected.
Choose five disjoint paths from $S$ to $N_G(v)$, stopped when they
first reach $N_G(v)$. They avoid $v$, and the paths at vertices
of $X_0$ are trivial.
Their five ends form a set containing $X_0$. Apply
Lemma~\ref{lem:star} to $H$ rooted at these ends. Its central branch
set avoids $S$. Use this branch set and $\{v\}$ as helpers. Four
linkage paths become root branch sets. For the fifth, take the path
ending in the central branch set and delete its last vertex. The
trimmed path is nonempty, since its original end lies outside $X_0$
and hence outside $S$.
Both helpers are adjacent to the other four root branch sets and to
each other, and the trimmed branch set is adjacent to the central
helper. Thus at least ten of the eleven helper adjacencies are present.
Every branch set lies in $A\cup S$, and no artificial boundary edge
has been used.

The opposite side is internally five-connected and has density
at least $m+7$, where $m=10-e(G[S])$, by
Lemma~\ref{gl:orientations}(2). It supplies a rooted clique on $S$.
Joining corresponding root branch sets gives a $\Kminus$ model,
a contradiction.
\end{proof}

The following elementary argument will exclude the remaining
five-cuts. A proper five-fragment is called an $F$-fragment, for a
family $F$ of edges, if its boundary contains both ends of an edge
in $F$. An $F$-end is an $F$-fragment containing no $F$-fragment as
a strict subset.

\begin{lemma}\label{gl:end}
Let $J$ be five-connected, with every proper five-fragment of order
at least four. If $A$ is an $F$-end, no edge $e\in F$ can meet $A$,
lie in $A\cup N_J(A)$ and have both ends in a five-cut of $J$.
\end{lemma}

\begin{proof}
Put $S=N_J(A)$ and $C=V(J)\setminus(A\cup S)$. Suppose a five-cut
$T$ contains both ends of such an edge $e$, and partition $J-T$
into two nonempty unions of components $B,D$. Write
\[
 \begin{array}{lll}
 p=|A\cap T|,&s=|S\cap T|,&r=|C\cap T|,\\
 u=|S\cap B|,&w=|S\cap D|.&
 \end{array}
\]
Then $p+s+r=u+s+w=5$, $p\geq1$ and $p+s\geq2$.

If $A\cap B\ne\varnothing$, its neighbourhood lies in
$(A\cap T)\cup(S\cap T)\cup(S\cap B)$. This set has order at least
five, by connectivity. If it has order five, it equals the boundary
and contains both ends of $e$, giving an $F$-fragment strictly
inside $A$. Its opposite side contains
$C$, and strict containment follows from $p\geq1$. Therefore
\[
 p+s+u\geq6,\qquad p>w,\qquad u>r.
\]
The opposite corner $C\cap D$ has boundary contained in a set of
order $r+s+w=5+r-u<5$, so it is empty. Symmetrically, if
$A\cap D\ne\varnothing$, then $p>u$, $w>r$ and $C\cap B$ is empty.
In each application, a nonempty opposite corner ensures that
the separation is proper.

If both $A$-corners are nonempty, then $C\subseteq T$, so
$r=|C|\geq4$, contrary to $p+s\geq2$. If just $A\cap B$ is
nonempty, then $D=S\cap D$, giving $w=|D|\geq4$. Since $p>w$,
we have $p=5$ and $r=s=0$. Now $C=C\cap B$ has all its neighbours
in $S\cap B$, of order $u=5-w\leq1$, again impossible. The other
one-corner case is symmetric.

If neither $A$-corner is nonempty, then $A\subseteq T$ and
$p=|A|\geq4$. Since $|C|\geq4$ and $r\leq1$, a $C$-corner is
nonempty, say $C\cap B$. Its boundary has order at most $r+s+u$,
so connectivity gives $u\geq p\geq4$. The other $C$-corner must
be empty, or similarly $w\geq4$. But then $D=S\cap D$, and again
$w=|D|\geq4$. This contradicts $u+w\leq5$.
\end{proof}

\begin{lemma}\label{gl:six-connected}
The graph $G$ is six-connected.
\end{lemma}

\begin{proof}
It is already five-connected. We first show that every proper
five-fragment has at least four vertices. A singleton would have
degree five. For a fragment of order $a=2$ or $3$, the degree bound
and the possible incident edges give
\[
 2a-\binom a2\leq\rho_G(A)\leq\binom a2+a.
\]
Its density is therefore three if $a=2$, and between three and six
if $a=3$. Both contradict Lemma~\ref{gl:orientations}(2), which
allows only densities at most one or at least seven. Also, every
low five-fragment is proper: an empty opposite side would give
density $18-e(G[N_G(A)])\geq8$.

Let $F$ be the edges with at most three common neighbours. Every
edge in $F$ has both ends in the boundary of a five-fragment of
density one. Indeed, its contraction preserves the density bound
by \eqref{gl:contraction-density}, so is not $4$-bilight. Its dense
bifragment lifts to two positive five-fragments, as described above;
Lemma~\ref{gl:orientations}(3) forces one of their densities to equal one.

Suppose a five-cut exists. It has a low side, so
Lemma~\ref{gl:low-edge} supplies an edge of $F$. The preceding
paragraph supplies a low $F$-fragment. Choose an inclusion-minimal
one, $A$. In fact $A$ is an $F$-end without any density restriction.
For a smaller $F$-fragment $B\subsetneq A$, a heavy $B$ would be
disjoint and nonadjacent to the heavy opposite side of $A$, violating
Lemma~\ref{gl:orientations}(3). Thus $B$ would be low, contradicting
the choice of $A$.

Lemma~\ref{gl:low-edge} now supplies an edge $e\in F$ meeting $A$.
Its other end lies in $A\cup N_G(A)$, and both ends lie in a
five-cut by the earlier contraction argument. This contradicts
Lemma~\ref{gl:end} and excludes every five-cut.
\end{proof}

\section{The final degree cases}
\label{sec:global-terminal}

The degree-six argument adapts the final part of the proof of
Norin and Totschnig in~\cite[Section~3]{NT2025}; six-connectivity
excludes the separator outcome directly. It uses
the two-paths theorem of Robertson, Seymour and
Thomas~\cite[Theorem~(2.4)]{RST1993}, in the form also
quoted in \cite[Theorem~13]{NT2025}. For four distinct vertices in
a specified cyclic order, it gives disjoint paths joining opposite
pairs, a separation of order at most three with all terminals on
one side and at least two vertices in the other open side, or a
drawing in a disc with the terminals on its boundary in that order.

\begin{lemma}\label{gl:degree-six-terminal}
Let $H$ be six-connected, with at least nine vertices and no
$\Kminus$ minor. If every edge has at least four common neighbours
and $H$ has a vertex of degree six, then $e(H)\leq4v(H)-9$.
\end{lemma}

\begin{proof}
Let $d_H(v)=6$. The complement of $H[N_H(v)]$ is a matching, since
each neighbour of $v$ has at least four neighbours there. Label the
six neighbours $u_1,w_1,u_2,w_2,u_3,w_3$ so that the only possible
nonedges are $u_iw_i$.

There cannot be disjoint paths joining two different pairs $u_i,w_i$
with interiors outside $N_H[v]$: contracting the interior of each
path into one of its ends supplies two of the three possibly missing
adjacencies. Together with $\{v\}$, the six resulting branch sets give
$\Kminus$. In particular all three pairs are nonedges. Otherwise
an edge joining one pair and a path for another pair through an exterior
component would suffice. Such a component exists and is adjacent
to all six neighbours of $v$, by six-connectivity.

Some pair, say $u_1,w_1$, has no common neighbour outside $N_H[v]$.
Otherwise choose an exterior common neighbour $x_i$ for each pair.
Two distinct choices give the forbidden disjoint length-two paths,
so all choices coincide at a vertex $x$. There is a component $C$
outside $N_H[v]\cup\{x\}$, because $v(H)\geq9$. Its boundary is
contained in $N_H(v)\cup\{x\}$ and has at least six vertices.
Thus at least five neighbours of $v$ meet $C$, giving a path for
one pair through $C$ and a disjoint path for another through $x$.

Apply the two-paths theorem to
$H'=H-\{v,u_1,w_1\}$, with ordered terminals $u_2,u_3,w_2,w_3$.
The two paths are excluded above. A separation $(A,B)$ in the
second outcome lifts to
\[
 (A\cup\{v,u_1,w_1\},\ B\cup\{u_1,w_1\})
\]
in $H$, of order at most five. It is proper: $v$ lies in the first
open side, and the second contains $B\setminus A$. No neighbour
of $v$ lies in that second open side. This contradicts
six-connectivity.

Hence $H'$ has the disc drawing from the third outcome. Its four
distinct boundary terminals give $e(H')\leq3v(H')-7$.
Every vertex outside $N_H[v]$
meets the deleted triple at most once, by the choice of $u_1,w_1$.
Each of the other four neighbours of $v$ meets it three times,
and the triple induces two edges. With $n'=v(H')=v(H)-3$,
\[
 e(H)\leq(3n'-7)+(n'-4)+12+2=4v(H)-9.
\]
\end{proof}

\begin{proof}[Proof of Theorem~\ref{thm:bilight}]
The minimum counterexample $G$ is six-connected by
Lemma~\ref{gl:six-connected}. Every edge has at least four common
neighbours: otherwise contracting it gives a five-connected, hence
$4$-bilight, smaller graph of density at least $-2$. The average
degree is $8-4/v(G)<8$, whereas Lemma~\ref{gl:degree-six} gives
minimum degree at least six. Thus some vertex has degree six or seven.

Degree six contradicts Lemma~\ref{gl:degree-six-terminal}, since
$e(G)=4v(G)-2$. For a degree-seven vertex $v$, put $J=G[N_G(v)]$.
The common-neighbour bound gives $\delta(J)\geq4$. Choose a
component $C$ outside $N_G[v]$, which exists since $v(G)\geq9$.
Its boundary lies in $N_G(v)$ and has at least six vertices.
Contract $C$ to a vertex and delete the other exterior components.
The resulting nine-vertex minor consists of $J$, the vertex $v$
adjacent to all of $J$, and a nonadjacent second vertex adjacent to
at least six vertices of $J$. Lemma~\ref{lem:quotient} gives a
$\Kminus$ model in this minor. Replacing the contracted vertex by
the connected set $C$ lifts the model to $G$.

Both degree cases are impossible, so no counterexample exists.
\end{proof}

\begin{proof}[Proofs of Corollary~\ref{thm:five-connected} and
Theorem~\ref{thm:colour}]
Every five-connected graph is $4$-bilight, giving
Corollary~\ref{thm:five-connected}. For the colouring theorem, suppose
a $\Kminus$-minor-free graph is not six-colourable, and choose such
a graph $G$ minor-minimal. Then its chromatic number is at least
seven and every proper minor is six-colourable.
The critical-graph reduction~\cite[Theorem~1.6]{DNR2026} gives
seven-connectivity and $e(G)\geq4v(G)-2$. The graph is therefore
$4$-bilight, and Theorem~\ref{thm:bilight} supplies the forbidden
minor.
\end{proof}


\appendix
\section{The finite degree-seven lemma}\label{sec:finite}

\begin{lemma}\label{lem:quotient}
Let $R$ be a graph on seven vertices with minimum degree at least four.
Add a vertex $v$ adjacent to every vertex of $R$, and another vertex
$c$ adjacent to at least six vertices of $R$. No edge $vc$ is required.
The resulting graph contains a $K_7^{-}$ minor.
\end{lemma}

This lemma is computer-assisted. Its verification concerns only
nine-vertex graphs; the global proof supplies the reduction from hosts
of arbitrary order. We describe the enumeration and the certificates
so that the scope of the computation is explicit.

The complement $F$ of $R$ has maximum degree at most two. Each component
of $F$ is therefore a path, an isolated vertex, or a cycle of length at
least three. Enumerating multisets of these components with total order
seven gives $29$ isomorphism types. For each type, the verifier tests all
eight possibilities for $c$: it meets every vertex of $R$, or misses
one specified vertex. Thus there are $29\cdot8=232$ cases. Extra edges,
including $vc$, can be discarded.

For each case, the program searches for seven bags among the nine
vertices. It enumerates every partition of a subset of order seven,
eight or nine into seven nonempty blocks. There are respectively
$36$, $252$ and $462$ such partitions. A certificate records the bags
as vertex bitmasks. The verifier checks that they are disjoint and
nonempty, that every bag is connected, and that at most one of the
$21$ pairs of bags lacks an edge. These checks certify a $K_7^{-}$
minor directly; no minor-testing library is used.

All $232$ cases pass. The returned models use seven vertices in
$67$ cases, eight in $102$ cases and nine in $63$ cases. The program
also checks a literal $K_7^{-}$ as a positive example and a padded
$K_6$ and $K_{2,2,2,2}$ as negative examples. It validates each returned
certificate after the search; this validation reuses the connectivity
and contact predicates and is not an independently implemented checker.

The complete standard-library Python program is preserved at a fixed
revision in the accompanying repository:
\begin{center}
\href{https://github.com/cavit99/hadwiger-k7/blob/320ea3275197f3b5124b38fb26af24bb05754d57/active/hc7_k7minus_degree7_quotient_verify.py}
{\texttt{hc7\_k7minus\_degree7\_quotient\_verify.py}}.
\end{center}
Run it with \texttt{uv run python3} from a checkout of that revision.
The manuscript's accompanying README records the source hash, expected
output and certificate digest.


\section*{Use of AI}
OpenAI models helped develop the proof and assisted with computation,
drafting and internal proof checks. The author guided the research plan.

\bibliographystyle{amsplain}
\begingroup
\raggedright
\bibliography{references}
\endgroup
\end{document}
